\begingroup%
\definecolor{ksBlue}{HTML}{3478B8}%
\definecolor{ksTeal}{HTML}{27877F}%
\definecolor{ksAmber}{HTML}{C58B28}%
\definecolor{ksInk}{HTML}{243447}%
\begin{tikzpicture}[
  x=1mm, y=1mm,
  card/.style={draw, fill=white, rounded corners=2mm, line width=.8pt,
    text=ksInk, align=flush center, font=\fontsize{9.5}{12.5}\selectfont,
    minimum height=20mm, inner xsep=2mm, inner ysep=1.7mm, outer sep=.4pt},
  flow/.style={-{Stealth[length=2.6mm,width=1.9mm]}, draw=ksInk, line width=.9pt},
  edge text/.style={text=ksInk, font=\fontsize{9}{12.5}\selectfont, inner sep=1mm}
]
  \node[card, draw=ksBlue, minimum width=100mm] (l1) at (0,0)
    {\textbf{第一层：Math\_reshape（数学变换）}\\
    等价改写计算顺序 · 在线状态递推 · 精度取舍\\
    难度最高，可省下的计算量与访存最多};
  \node[card, draw=ksTeal, minimum width=120mm] (l2) at (0,-33)
    {\textbf{第二层：Work\_schedule（Tile 以上的全局调度）}\\
    同构 Worktile 切分、SM / CTA / Cluster 映射与负载均衡\\
    难度中等，决定设备级并行利用率};
  \node[card, draw=ksAmber, minimum width=140mm] (l3) at (0,-66)
    {\textbf{第三层：Tile\_pipeline（Tile 内微流水线）}\\
    Microbench 周期先验 · Warp 角色划分 · 8 动词与等待边编排异步重叠\\
    问题最具体，决定能多接近硬件峰值吞吐};

  \draw[flow] (l1) -- node[edge text, right=1mm]
    {确定计算哪些量、跨轮携带哪些状态} (l2);
  \draw[flow] (l2) -- node[edge text, right=1mm]
    {确定单个 Worker 的子任务边界} (l3);
\end{tikzpicture}%
\endgroup%
